\section{Introduction}
\label{sec:introduction}

The estimation of a time-varying conditional quantile from a single observed time series is a fundamental problem in nonparametric statistics, arising in environmental monitoring, biomedical signal processing, finance, and astrophysics, among other fields. Unlike mean regression, quantile estimation provides a direct view of the upper or lower tails of the response distribution, which are often of greater scientific interest than the central tendency. When the underlying trend is believed to be piecewise polynomial---constant level shifts, piecewise linear behavior, or piecewise quadratic curvature---trend filtering provides a flexible and computationally tractable framework, producing estimates that adapt automatically to abrupt changes while remaining smooth in regions of slow variation.

The frequentist approach to quantile trend filtering was developed by \citet{brantley2020baseline}, who combined the asymmetric Laplace check loss of \citet{koenker1978regression} with the discrete difference penalty of \citet{tibshirani2014adaptive} and solved the resulting nonsmooth convex program via the alternating direction method of multipliers (ADMM). Their estimator is fast, scales well to large signals, and inherits the locally adaptive properties of $\ell_1$ trend filtering. Theoretical properties of the resulting estimator, including minimax risk bounds over piecewise polynomial signal classes, have been established by \citet{madrid2022risk}. The procedure is, however, an optimization solver: it returns a single point estimate of the trend at a prespecified penalty level, with no associated measure of uncertainty. Credible or confidence intervals must be obtained externally---typically through the bootstrap, which does not address the regularization parameter selection problem and inflates the computational burden by a factor equal to the number of resamples. The scale of the asymmetric Laplace likelihood is, without loss of generality, fixed at unity in the ADMM formulation, since it does not affect the location of the optimum; this convention disappears as soon as one wishes to treat the model as a probabilistic object.

These limitations motivate a fully Bayesian treatment of the same problem. A Bayesian formulation of quantile trend filtering yields a posterior distribution over the trend, providing credible intervals and other coherent uncertainty summaries as a direct byproduct of inference, while allowing the regularization parameter to be inferred from the data rather than tuned externally. Related progress on Bayesian quantile regression in the additive and spline settings has been made by \citet{fasiolo2021fast}, though their approach targets smooth additive models rather than piecewise polynomial trends of arbitrary order. The analogous trajectory for mean trend filtering has been traced by Heng, Zhou, and Chi, who placed the $\ell_1$ trend filtering problem of \citet{tibshirani2014adaptive} in a Bayesian framework using proximal Markov chain Monte Carlo techniques \citep{pereyra2016proximal} to handle the nondifferentiable prior and obtain posterior samples along with automatic, data-driven calibration of the smoothing parameter. The quantile setting inherits the same motivation but introduces a substantive additional difficulty: the asymmetric Laplace likelihood is itself nonsmooth, so a Bayesian quantile trend filtering sampler must contend with two sources of nondifferentiability rather than one.

A closely related Bayesian quantile trend filtering procedure is the Gibbs sampler of \citet{onizuka2024fast}, which exploits the scale-mixture representation of the asymmetric Laplace distribution to construct conjugate full conditionals under either a Bayesian Lasso or a horseshoe prior placed locally on each coordinate of the difference vector $D\theta$. The resulting sampler is principled, provides genuine posterior inference, and achieves strong empirical performance, particularly on piecewise constant signals where the per-coefficient horseshoe prior delivers high local adaptivity. Two structural limitations, however, restrict its practical reach. First, the trend update requires sampling a multivariate Gaussian whose precision matrix is $D^\top \mathrm{diag}(w)^{-1} D + \mathrm{diag}(c)$, where the diagonal weights $w$ are themselves random variables resampled at every Gibbs iteration. The matrix therefore changes from one iteration to the next, foreclosing any possibility of caching its factorization, and its factorization must be recomputed at every step, which in the reference implementation is done with dense linear algebra, making the exact MCMC sampler slow even at moderate sample sizes. Second, the variational Bayes (VB) approximation proposed as a faster alternative relies on ratios of modified Bessel functions of high order, which become numerically unstable at moderate sample sizes; and even were these numerical issues resolved, mean-field variational approximations are well known to systematically underestimate posterior variance, undermining the very inferential guarantees that motivate the Bayesian framework.

The present paper develops a Bayesian quantile trend filtering procedure that retains the scalability and structural advantages of the frequentist ADMM solver while delivering full posterior inference. The model is the standard one: an asymmetric Laplace likelihood at the target quantile level $\tau$ and a global Laplace prior on the $(k+1)$-th order differences of the trend, equivalent in posterior mode to the \citet{brantley2020baseline} estimator. The contribution lies in the sampler. We construct a three-block Gibbs scheme in which the regularization parameter $\lambda$ is updated from a conjugate Gamma full conditional, the asymmetric Laplace scale $\sigma$---treated as an unknown parameter rather than fixed at unity, in contrast to the optimization formulation---is updated from a conjugate inverse-Gamma full conditional, and the trend $\theta$ is updated via a single step of proximal Metropolis-adjusted Langevin algorithm (P-MALA) \citep{pereyra2016proximal}. The proximal map required by P-MALA is itself a constrained $\ell_1$ trend filtering problem, which we solve by a small number of ADMM iterations whose linear system has a fixed coefficient matrix that depends only on the deterministic step sizes and the difference operator. Because this matrix does not change between Gibbs iterations and is sparse and banded, it admits a Cholesky factorization that is computed once and reused thereafter, with backsolve cost linear in $n$. The sampler is therefore fully Bayesian, provides automatic and self-calibrating updates for both the regularization strength and the noise scale, and exhibits per-iteration complexity that is substantially lower than the dense Gibbs sampler of \citet{onizuka2024fast}.

The framework extends naturally to joint inference across multiple quantile levels under the standard non-crossing constraint. The single-quantile P-MALA step is replaced by a sequential update in which the proximal map is solved under box constraints determined by the most recent values of the neighboring quantile trends. The required modification to the ADMM solver introduces a single auxiliary variable and an elementwise projection step, leaving its overall structure and computational profile intact. This joint construction, however, also exposes a difficulty intrinsic to gradient-based samplers under the asymmetric Laplace likelihood at extreme quantile levels. For $\tau$ close to zero or one, nearly all residuals share the same sign and the gradient of the smoothed check loss approaches a constant vector, so the P-MALA proposal carries negligible directional information and the chain develops a persistent drift away from the data. The locally adaptive prior of \citet{onizuka2024fast} partially mitigates this pathology, and their sampler produces degraded but not catastrophic estimates at extreme quantile levels. We instead resolve it through a strategy that exploits the box-constraint structure of the joint sampler: we augment the set of inferential quantile levels with auxiliary \emph{sentinel} quantiles placed beyond the target range, whose role is purely structural. The sentinels are sampled as ordinary members of the joint posterior and provide a dynamically updating upper or lower bound on the extreme target quantile through the non-crossing constraint, confining its exploration to a statistically sensible region at every iteration. The mechanism is computationally cheap and, in our simulation study, allows the proposed sampler to match or slightly improve upon the extreme-quantile accuracy of \citet{onizuka2024fast} while using a globally penalized prior that is structurally simpler than the per-coefficient horseshoe.

The main contributions of this work are as follows.
\begin{enumerate}
    \item We introduce a Bayesian quantile trend filtering sampler in which the trend, the regularization parameter, and the asymmetric Laplace scale are all inferred jointly from conjugate or nearly conjugate full conditionals, embedding a proximal MALA step for the nondifferentiable component of the posterior. The scale parameter is treated as an unknown, breaking with the optimization convention of fixing it at unity and allowing the effective likelihood temperature to adapt to the data.
    \item We exploit the deterministic structure of the proximal subproblem to obtain a sampler whose per-iteration cost is linear in the sample size, in contrast to the cubic per-iteration cost of the existing Bayesian competitor. The coefficient matrix of the inner ADMM linear system is sparse, banded, and constant within and across Gibbs iterations, so its factorization is computed once and reused, eliminating the dominant computational bottleneck of the locally adaptive Gibbs approach of \citet{onizuka2024fast}.
    \item We extend the sampler to the joint estimation of multiple quantile levels under non-crossing constraints, incorporating the constraints exactly through a box-projection step inside the inner ADMM solver rather than through approximate or penalty-based enforcement.
    \item We introduce the sentinel quantile strategy, a structurally minimal device that resolves the gradient degeneracy of proximal MALA at extreme quantile levels and recovers competitive accuracy at the tails of the response distribution. The mechanism is generic and applies to any gradient-based sampler for quantile trend filtering, not only the one developed here.
    \item Through a comprehensive simulation study spanning four signal structures, three noise distributions, three difference orders, and both interior and extreme quantile levels, we demonstrate that the proposed sampler delivers accuracy comparable to or better than the competitive Bayesian approach of \citet{onizuka2024fast}, with substantially reduced computational cost and proper posterior uncertainty quantification, and that it does so using a single global penalty parameter rather than the per-coordinate local shrinkage employed by competing priors, which simplifies the prior specification and the resulting sampler.
\end{enumerate}

The remainder of the paper is organized as follows. Section~\ref{sec:methodology} develops the methodology, presenting the data model and likelihood, the individual-quantile sampler with its three conjugate and proximal updates, the joint sampler with non-crossing constraints and the associated ADMM proximal operator, and the sentinel quantile strategy. Section~\ref{sec:simulations} reports the simulation study comparing the proposed sampler with the relevant frequentist and Bayesian competitors across interior and extreme quantile levels. Section~\ref{sec:discussion} discusses the results, limitations, and avenues for extension.